\begin{pastexam}[label={px:ntu-112-dsa-16}]{臺大資工所 112 學年度 第 16 題}
（問答；只看最後答案，完全正確才給分）

\textbf{(9 points)} The pseudo code function to compute the prefix function in the
Knuth-Morris-Pratt (KMP) string matching algorithm is given below. Given a pattern string
$P[1..m]$, the prefix function for this pattern $P$ is the function
$\pi:\{1,2,\ldots,m\}\to\{0,1,\ldots,m-1\}$ such that $\pi[q]$ is the length of the longest prefix
of P that is a proper suffix of $P_q$, where $P_q$ denotes the $q$-character prefix of the string
$P$. For pattern string $P$=ABACABACABACABAD, how many times is \textbf{line 7} executed? (Note:
only the final answer will be graded and only fully correct answer will be given points)

\smallskip
\textsc{Compute-Prefix-Function}(P)
\begin{tabbing}
\hspace{2em}\=\hspace{2em}\=\hspace{2em}\=\hspace{2em}\=\kill
\llap{1}\>$m = P.\mathit{length}$\\
\llap{2}\>let $\pi[1..m]$ be a new array\\
\llap{3}\>$\pi[1] = 0$\\
\llap{4}\>$k = 0$\\
\llap{5}\>\textbf{for} $q = 2$ \textbf{to} $m$\\
\llap{6}\>\>\textbf{while} $k > 0$ and $P[k+1] \ne P[q]$\\
\llap{7}\>\>\>$k = \pi[k]$\\
\llap{8}\>\>\textbf{if} $P[k+1] == P[q]$\\
\llap{9}\>\>\>$k = k+1$\\
\llap{10}\>\>$\pi[q] = k$\\
\llap{11}\>\textbf{return} $\pi$
\end{tabbing}
\end{pastexam}

\begin{pastsoln}
答案是 $5$ 次。

題目的第 7 行是失配時退到較短邊界的 $k = \pi[k]$，就是資結書虛擬碼~27.3 的第 6 行。
$P = \mathtt{ABACABACABACABAD}$，$m = 16$。 逐一列出每個 $q$ 開始時的 $k$、第 7 行執行的情形和
$\pi[q]$：

\begin{center}
\small
\begin{tabular}{ccc>{\raggedright\arraybackslash}p{0.5\linewidth}cc}
\toprule
$q$ & $P[q]$ & 開始的 $k$ & 題目第 6 行的比較與第 7 行 & 次數 & $\pi[q]$ \\
\midrule
2 & \texttt{B} & 0 & $k = 0$，不進迴圈；$P[1] = \mathtt{A} \ne \mathtt{B}$ & 0 & 0 \\
3 & \texttt{A} & 0 & 不進迴圈；$P[1] = \mathtt{A}$ 相同 & 0 & 1 \\
4 & \texttt{C} & 1 & $P[2] = \mathtt{B} \ne \mathtt{C}$，$k \gets \pi[1] = 0$ & 1 & 0 \\
5 & \texttt{A} & 0 & 不進迴圈；$P[1]$ 相同 & 0 & 1 \\
6--15 & & $q - 5$ & 每次 $P[k+1] = P[q]$，迴圈條件一開始就不成立 & 0 & $q - 4$ \\
16 & \texttt{D} & 11 & $P[12] = \mathtt{C}$，$k \gets \pi[11] = 7$；$P[8] = \mathtt{C}$，$k \gets \pi[7] = 3$；
$P[4] = \mathtt{C}$，$k \gets \pi[3] = 1$；$P[2] = \mathtt{B}$，$k \gets \pi[1] = 0$ & 4 & 0 \\
\bottomrule
\end{tabular}
\end{center}

$q = 6$ 到 $15$：$P[5..15]$ 和 $P[1..11]$ 相同（\texttt{ABAC} 一直重複），每一步 $P[k + 1] = P[q]$
都直接接上，$k$ 從 $1$ 一路加到 $11$。 $q = 16$ 的 \texttt{D} 不在前面出現過，$k$ 沿著邊界的鏈
$11 \to 7 \to 3 \to 1 \to 0$ 退到底，每一步比的 $P[k + 1]$ 是 \texttt{C}、\texttt{C}、\texttt{C}、
\texttt{B}，都不是 \texttt{D}。 合計 $1 + 4 = 5$ 次。 完整的
$\pi = 0, 0, 1, 0, 1, 2, 3, 4, 5, 6, 7, 8, 9, 10, 11, 0$。

檢查：資結書定理~27.4 的位能論證說總次數不超過 $m - 1 = 15$，而且 $k$ 一共只加了
$1 + 1 + 10 = 12$ 次（$q = 3$、$5$ 和 $6$ 到 $15$ 各一次），退的次數不會超過加的次數，$5 \le 12$。
\end{pastsoln}
